\section{Examples of Fields}

The simplest example is a scalar field
\[
    \phi : \mathbb{R}\times\mathbb{R}^3 \to \mathbb{R},
\]
where \(\phi(t,\vect{x})\) assigns one real number to each point of space at
each time. Temperature in a material body is a familiar nonrelativistic scalar
field.

A vector field assigns a vector to each point. The velocity field of a fluid is
written
\[
    \vect{v}(t,\vect{x}) =
    \bigl(v^1(t,\vect{x}),v^2(t,\vect{x}),v^3(t,\vect{x})\bigr).
\]
It describes the local velocity of the material element passing through
\(\vect{x}\) at time \(t\).

Electromagnetism is naturally expressed in terms of fields. In elementary
vector notation one uses the electric and magnetic fields
\[
    \vect{E}(t,\vect{x}), \qquad \vect{B}(t,\vect{x}).
\]
In relativistic notation these are combined into the electromagnetic field
strength tensor \(F_{\mu\nu}\).

\begin{examplebox}
At a fixed time, the temperature inside a room is not one number. It is a
function \(T(\vect{x})\). Heating, cooling, and diffusion are statements about
how this function changes in time and space.
\end{examplebox}
